\section{From Sources to Date Constraints}
\label{sec:sources}

Temporal evidence contains several kinds of dates.
The retrieval interface must preserve their meaning before applying any replacement rule.
We distinguish an uncertain occurrence from a known duration and from a dated observation.
Table~\ref{tab:date-types} gives the operational difference.

\begin{table}[t]
\centering\small
\begin{tabular}{@{}p{.28\columnwidth}p{.65\columnwidth}@{}}
\toprule
Source meaning & Representation and operation\\
\midrule
Update occurred in a year & Bound its unknown start by that year's dates; use the two-witness compiler.\\[3pt]
Fact held during a stated period & Preserve its validity interval; test overlap with the question's period.\\[3pt]
Fact was observed in a yearly snapshot & Preserve the observation and all listed answers; do not infer an exact transition date.\\
\bottomrule
\end{tabular}
\caption{Three date meanings require different operations. The compiler takes bounds on one occurrence date.}
\label{tab:date-types}
\end{table}

\paragraph{Occurrence precision.}
An exact day gives a point on the chosen time axis.
A stated month gives bounds on an occurrence within that month.
A stated year gives wider bounds.
The representation keeps this precision instead of substituting January 1 or the interval midpoint.
Relative expressions require a source anchor.
An expression such as \emph{around 1713} does not itself supply a numerical error radius.
The source adapter must provide justified bounds before invoking the compiler.

\paragraph{Replacement decisions.}
A directed decision states that one claim replaces another if it occurs strictly later.
The decision concerns the same subject and attribute.
It must distinguish a changed value from an additional valid value.
For example, the Chicago Stadium passage gives overlapping periods for the Blackhawks and the Bulls.
Neither tenancy replaces the other.
The compiler takes pair decisions as inputs, together with source identifiers and date bounds.
This interface allows the same executor to serve rule-based and learned extractors.

\paragraph{Question scope.}
We study questions asking which claims apply at some time within a specified period.
Point questions are the special case with identical start and end dates.
A question asking what held throughout a period requires a different predicate.
The query adapter preserves this distinction.
It preserves the supplied date precision.
Thus a question about a whole year does not silently become a question about its last day.

\paragraph{Source units.}
Replacement operates on claims, rather than entire paragraphs.
A single paragraph can contain several historical states and several relations.
The News Quiz passage, for example, lists successive hosts and a returning host.
Removing that paragraph because one host was replaced would remove other requested evidence.
Each extracted claim retains its passage identifier and source span.
The retrieval layer can recover the supporting text after temporal selection.

\paragraph{Three output states.}
For a given question, the compiler marks each claim as guaranteed, unresolved, or unsupported.
A guaranteed claim has support under every permitted timeline.
An unresolved claim has support under some, but not all, timelines.
An unsupported claim has support under none.
The interface exposes these states before the context budget is applied.
This preserves the difference between lacking temporal support and merely having uncertain dates.
